\begin{abstract}
Adaptive question selection in educational systems is fundamentally limited by \emph{forgetting}: knowledge decays over time for topics that are not reinforced, rendering static selection policies suboptimal.
We present a comprehensive study of 14 bandit algorithms for adaptive question selection under Ebbinghaus forgetting dynamics, framing the problem as a restless multi-armed bandit (RMAB).
We introduce four novel algorithms---Forgetting-UCB (F-UCB), which augments exploration bonuses with decay-aware urgency; BKT-Bandit, which maintains Bayesian posteriors with integrated forgetting; Whittle-advantage, which computes budget-aware indices via MDP value differences; and MetaSelector, a UCB-over-experts approach that interpolates across regimes online.
Ours is the first work to formalize educational quizzing as a restless bandit and to identify the \emph{Oracle catastrophe}: a clairvoyant selector that exploits current knowledge states loses to random selection in 10 out of 12 configurations due to systematic neglect of forgetting arms.
Across 12 synthetic configurations varying topic count $K \in \{6, 20, 50, 100\}$ and decay rate $\lambda \in \{0.005, 0.01, 0.05\}$, BKT-Bandit achieves top-3 performance in 10/12 settings with $+8.2\%$ average improvement over random, F-UCB delivers up to $+41.8\%$ improvement at high decay, and MetaSelector ranks in the top-5 in 11/12 configurations with principled regime interpolation.
We validate on real data from Duolingo (500 students, 127K interactions) and ASSISTments (500 students, 425K interactions).
Code and data are publicly available.
\end{abstract}
